\documentclass[11pt,a4paper]{article}

\usepackage[T1]{fontenc}
\usepackage[utf8]{inputenc}
\usepackage[ngerman]{babel}

\usepackage{amsmath}
\usepackage{amssymb}
\usepackage{amsthm}
\usepackage{mathtools}
\usepackage{geometry}
\usepackage{hyperref}
\usepackage{graphicx}
\usepackage{enumitem}

\geometry{margin=2.5cm}

\DeclareGraphicsExtensions{.pdf,.png,.jpg}
\graphicspath{{./}{fig/}{figs/}{images/}{img/}}

\newtheorem{satz}{Satz}[section]
\newtheorem{lemma}[satz]{Lemma}
\newtheorem{korollar}[satz]{Korollar}

\theoremstyle{definition}
\newtheorem{definition}[satz]{Definition}
\newtheorem{beispiel}[satz]{Beispiel}

\theoremstyle{remark}
\newtheorem{bemerkung}[satz]{Bemerkung}

\title{
	Bamberger Modell III\\
	Eine tetraedrische Struktur auf den natürlichen Zahlen\\
	Quaternionische Orientierung, modulare Phasen und ikosaedrische Geometrie
}

\author{
	Thomas Hoffbauer
}

\date{\today}

\begin{document}
	
	\maketitle
	
	\begin{abstract}
		
		Diese Arbeit entwickelt eine geometrische Sichtweise auf die natürlichen Zahlen.
		Ausgehend von den invertierbaren Restklassen modulo \(12\) wird eine vierstufige
		Struktur konstruiert,
		
		\[
		V_4
		\rightarrow
		Q_8
		\rightarrow
		\eta_{24}
		\rightarrow
		A_5,
		\]
		
		welche Restklassenarithmetik, Quaternionen,
		modulare Formen und Ikosaedersymmetrien miteinander verbindet.
		
		Im Zentrum steht die Beobachtung, dass die vier Klassen
		
		\[
		1,5,7,11 \pmod{12}
		\]
		
		eine natürliche tetraedrische Geometrie erzeugen.
		Durch eine quaternionische Hebung entsteht eine orientierte Spinorstruktur.
		Die Dedekindsche Eta-Funktion liefert eine diskrete
		24-periodische Phasenstruktur,
		während die Ikosaedergruppe \(A_5\)
		eine natürliche Bahngeometrie bereitstellt.
		
		Die Arbeit formuliert einen geometrischen Rahmen,
		in dem natürliche Zahlen durch Skala,
		Spin,
		Bindung,
		Orbit
		und Phase beschrieben werden können.
		
	\end{abstract}
	
	\tableofcontents
	
	\newpage
	
	%%%%%%%%%%%%%%%%%%%%%%%%%%%%%%%%%%%%%%%%%%%%%%%%%%%%%%%%%%%%
	
	\section{Einleitung}
	
	Die klassische Zahlentheorie beschreibt natürliche Zahlen
	über ihre Größe,
	ihre Teilbarkeit
	und ihre Zerlegung in Primfaktoren.
	
	Seit Euklid gilt:
	
	\begin{quote}
		Jede natürliche Zahl zerfällt eindeutig in Primzahlen.
	\end{quote}
	
	Diese Tatsache bildet die Grundlage großer Teile der modernen Mathematik.
	
	Das vorliegende Modell geht von einer weiterführenden Frage aus:
	
	\begin{quote}
		Besitzen Primfaktorisierungen neben ihrer arithmetischen
		auch eine geometrische Struktur?
	\end{quote}
	
	Die Arbeit untersucht die Möglichkeit,
	dass Primzahlen nicht nur als Zahlen,
	sondern zugleich als Träger einer verborgenen
	Orientierungsstruktur verstanden werden können.
	
	Dabei entsteht die Hierarchie
	
	\[
	V_4
	\rightarrow
	Q_8
	\rightarrow
	\eta_{24}
	\rightarrow
	A_5.
	\]
	
	Diese verbindet vier klassische Gebiete der Mathematik:
	
	\begin{itemize}
		\item Restklassenarithmetik,
		\item Quaternionenalgebra,
		\item modulare Formen,
		\item Polyedersymmetrien.
	\end{itemize}
	
	%%%%%%%%%%%%%%%%%%%%%%%%%%%%%%%%%%%%%%%%%%%%%%%%%%%%%%%%%%%%
	
	\section{Historische Wurzeln}
	
	\subsection{Euler und die Primzahlen}
	
	Die moderne Zahlentheorie beginnt mit Euler.
	
	Die Euler-Produktdarstellung
	
	\[
	\zeta(s)
	=
	\prod_p
	\frac{1}{1-p^{-s}}
	\]
	
	macht deutlich,
	dass Primzahlen die elementaren Bausteine
	der natürlichen Zahlen darstellen.
	
	Alle natürlichen Zahlen entstehen aus deren Multiplikation.
	
	Im Bamberger Modell werden Primzahlen
	als elementare Orientierungsträger betrachtet.
	
	\subsection{Hamilton und die Quaternionen}
	
	1843 entdeckte Hamilton die Quaternionenalgebra
	
	\[
	\mathbb H
	=
	\langle 1,i,j,k\rangle.
	\]
	
	Sie erfüllt
	
	\[
	i^2=j^2=k^2=ijk=-1.
	\]
	
	Zum ersten Mal erscheint eine Zahlstruktur,
	in der die Multiplikation nicht kommutativ ist.
	
	Quaternionen bilden heute die Grundlage
	von Rotationen und Spinoren.
	
	\subsection{Hurwitz}
	
	Hurwitz zeigte,
	dass die Quaternionen 24 ausgezeichnete Einheiten besitzen.
	
	Diese bilden den berühmten 24-Zeller.
	
	Die Zahl
	
	\[
	24
	\]
	
	wird später in der Theorie
	über die Eta-Funktion erneut auftreten.
	
	\subsection{Dedekind}
	
	Dedekind definierte die Eta-Funktion
	
	\[
	\eta(\tau)
	=
	q^{1/24}
	\prod_{n\ge1}(1-q^n).
	\]
	
	Diese Funktion spielt eine fundamentale Rolle
	in der Theorie modularer Formen.
	
	\subsection{Ramanujan}
	
	Ramanujan entdeckte tiefe Beziehungen zwischen
	
	\[
	e^{-\pi},
	\qquad
	e^{-2\pi},
	\qquad
	\varphi
	=
	\frac{1+\sqrt5}{2}.
	\]
	
	Insbesondere in den Rogers--Ramanujan-Identitäten
	erscheinen diese Größen gemeinsam.
	
	\subsection{Klein}
	
	Felix Klein erkannte,
	dass die Ikosaedergruppe
	
	\[
	A_5
	\]
	
	eng mit modularen Funktionen verbunden ist.
	
	Das Ikosaeder stellt die höchste reguläre
	Rotationssymmetrie des dreidimensionalen Raumes dar.
	
	%%%%%%%%%%%%%%%%%%%%%%%%%%%%%%%%%%%%%%%%%%%%%%%%%%%%%%%%%%%%
	
	\section{Die EABC-Algebra}
	
	\begin{definition}
		
		Die invertierbaren Restklassen modulo \(12\)
		
		\[
		(\mathbb Z/12\mathbb Z)^\times
		=
		\{1,5,7,11\}
		\]
		
		werden bezeichnet durch
		
		\[
		E,A,B,C.
		\]
		
	\end{definition}
	
	Es gilt
	
	\[
	E=1,
	\qquad
	A=5,
	\qquad
	B=7,
	\qquad
	C=11.
	\]
	
	Ferner:
	
	\[
	A^2=B^2=C^2=E.
	\]
	
	Außerdem:
	
	\[
	AB=C,
	\]
	
	\[
	BC=A,
	\]
	
	\[
	CA=B.
	\]
	
	\begin{satz}
		
		Die Elemente
		
		\[
		E,A,B,C
		\]
		
		bilden die Klein-Vierergruppe
		
		\[
		V_4.
		\]
		
	\end{satz}
	
	%%%%%%%%%%%%%%%%%%%%%%%%%%%%%%%%%%%%%%%%%%%%%%%%%%%%%%%%%%%%
	
	\section{Das tetraedrische Spinmodell}
	
	Wir definieren
	
	\[
	E=(1,1,1),
	\]
	
	\[
	A=(-1,-1,1),
	\]
	
	\[
	B=(1,-1,-1),
	\]
	
	\[
	C=(-1,1,-1).
	\]
	
	\begin{satz}
		
		Die vier Zustände bilden die Ecken eines regulären Tetraeders.
		
	\end{satz}
	
	\begin{proof}
		
		Zwischen jedem Paar ergibt sich die Distanz
		
		\[
		2\sqrt2.
		\]
		
		Alle sechs Kanten sind gleich lang.
		
	\end{proof}
	
	\begin{korollar}
		
		Es gilt
		
		\[
		E+A+B+C=0.
		\]
		
	\end{korollar}
	
	Dies zeigt,
	dass der vollständige Vierer spinneutral ist.
	
	%%%%%%%%%%%%%%%%%%%%%%%%%%%%%%%%%%%%%%%%%%%%%%%%%%%%%%%%%%%%
	
	\section{Kern-Hülle-Struktur}
	
	Die numerischen Untersuchungen legen folgende Interpretation nahe:
	
	\[
	E=\text{Kern},
	\]
	
	\[
	B,C=\text{Hülle},
	\]
	
	\[
	A=\text{Bindung}.
	\]
	
	Der elementare Hüllendipol lautet
	
	\[
	D_{BC}=B+C.
	\]
	
	Einsetzen ergibt
	
	\[
	D_{BC}
	=
	(0,0,-2).
	\]
	
	Die Gegenorientierung lautet
	
	\[
	-D_{BC}
	=
	(0,0,2).
	\]
	
	Damit besitzt die Hülle zwei natürliche Orientierungen.
	
	%%%%%%%%%%%%%%%%%%%%%%%%%%%%%%%%%%%%%%%%%%%%%%%%%%%%%%%%%%%%
	
	\section{Quaternionische Hebung}
	
	Die EABC-Algebra ist kommutativ.
	
	Zur Einführung einer Orientierung wird die Zuordnung
	
	\[
	E\mapsto1,
	\]
	
	\[
	A\mapsto i,
	\]
	
	\[
	B\mapsto j,
	\]
	
	\[
	C\mapsto k
	\]
	
	eingeführt.
	
	Dann gilt
	
	\[
	jk=i,
	\]
	
	aber
	
	\[
	kj=-i.
	\]
	
	Daraus folgt
	
	\[
	BC=+A,
	\]
	
	\[
	CB=-A.
	\]
	
	\begin{satz}
		
		Die quaternionische Hebung induziert
		eine orientierte Spinorstruktur.
		
	\end{satz}
	
	Die Orientierung entsteht nicht in \(V_4\),
	sondern erst in \(Q_8\).
	
	%%%%%%%%%%%%%%%%%%%%%%%%%%%%%%%%%%%%%%%%%%%%%%%%%%%%%%%%%%%%
	
	\section{Eta-24-Phasen}
	
	Die Dedekindsche Eta-Funktion erfüllt
	
	\[
	\eta(\tau+1)
	=
	e^{\pi i/12}\eta(\tau).
	\]
	
	Dies motiviert die Einführung einer Phase
	
	\[
	\varepsilon
	\in
	\mathbb Z_{24}.
	\]
	
	Wir setzen
	
	\[
	+A
	\leftrightarrow
	0,
	\]
	
	\[
	-A
	\leftrightarrow
	12.
	\]
	
	Die Vertauschung
	
	\[
	BC
	\leftrightarrow
	CB
	\]
	
	entspricht einem Halbumlauf
	
	\[
	\varepsilon
	\rightarrow
	\varepsilon+12.
	\]
	
	Nach zwei Vertauschungen
	
	\[
	12+12=24
	\]
	
	wird die Ausgangsphase wieder erreicht.
	
	%%%%%%%%%%%%%%%%%%%%%%%%%%%%%%%%%%%%%%%%%%%%%%%%%%%%%%%%%%%%
	
	\section{Ramanujan und der Goldene Schnitt}
	
	Ramanujans Arbeiten verbinden die Zahlen
	
	\[
	e^{-2\pi}
	\]
	
	und
	
	\[
	\varphi.
	\]
	
	Im Modell entsteht die Kette
	
	\[
	\eta_{24}
	\longrightarrow
	e^{-2\pi}
	\longrightarrow
	\varphi.
	\]
	
	Der Goldene Schnitt erscheint somit
	als geometrische Manifestation
	einer modularen Struktur.
	
	%%%%%%%%%%%%%%%%%%%%%%%%%%%%%%%%%%%%%%%%%%%%%%%%%%%%%%%%%%%%
	
	\section{Die ikosaedrische Erweiterung}
	
	Die Ikosaedergruppe
	
	\[
	A_5
	\]
	
	stellt die höchste reguläre Rotationssymmetrie
	des dreidimensionalen Raumes dar.
	
	Ihre Geometrie enthält den Goldenen Schnitt intrinsisch.
	
	Die ikosaedrische Ebene wird als orbitale Erweiterung
	des tetraedrischen Spinraums interpretiert.
	
	%%%%%%%%%%%%%%%%%%%%%%%%%%%%%%%%%%%%%%%%%%%%%%%%%%%%%%%%%%%%
	
	\section{Radix-Geometrie}
	
	Jede natürliche Zahl wird durch
	
	\[
	n=(G,H,V,\varepsilon)
	\]
	
	beschrieben.
	
	Dabei bezeichnet
	
	\[
	G
	\]
	
	den glatten Anteil,
	
	\[
	H
	\]
	
	den Spinanteil,
	
	\[
	V
	\]
	
	die Bindungs- und Orbitstruktur,
	
	und
	
	\[
	\varepsilon
	\]
	
	die Eta-Phase.
	
	Diese vier Größen bilden gemeinsam
	den geometrischen Zustand einer Zahl.
	
	%%%%%%%%%%%%%%%%%%%%%%%%%%%%%%%%%%%%%%%%%%%%%%%%%%%%%%%%%%%%
	
	\section{Arithmetische Teilchenspektroskopie}
	
	Die folgenden Begriffe sind mathematische Definitionen
	und keine physikalischen Identifikationen.
	
	\subsection{Spinoren}
	
	Eine einzelne Primzahl
	
	\[
	p>3
	\]
	
	wird als elementarer Spinor bezeichnet.
	
	\subsection{Mesonen}
	
	Zweierkombinationen
	
	\[
	A+B,
	\quad
	A+C,
	\quad
	B+C
	\]
	
	werden als Mesonen bezeichnet.
	
	\subsection{Baryonen}
	
	Dreierkombinationen
	
	\[
	A+B+C
	\]
	
	werden als Baryonen bezeichnet.
	
	\subsection{Singuletts}
	
	Der vollständige Vierer
	
	\[
	E+A+B+C=0
	\]
	
	bildet ein Singulett.
	
	%%%%%%%%%%%%%%%%%%%%%%%%%%%%%%%%%%%%%%%%%%%%%%%%%%%%%%%%%%%%
	
	\section{Masse als geometrische Rigidität}
	
	Im Modell wird Masse nicht als Substanzmenge verstanden.
	
	Masse beschreibt die Stabilität
	einer gekoppelten Struktur.
	
	Formal schreiben wir
	
	\[
	m
	=
	\Lambda(\eta,\varphi)
	\left(
	\Omega
	+
	\kappa \|S\|^2
	+
	\lambda \|L\|^2
	\right).
	\]
	
	Hierbei bezeichnet
	
	\[
	\Omega
	\]
	
	eine Faktorisierungstiefe,
	
	\[
	S
	\]
	
	die Spinstruktur,
	
	und
	
	\[
	L
	\]
	
	eine orbitale Struktur.
	
	Diese Formel besitzt derzeit heuristischen Charakter.
	
	%%%%%%%%%%%%%%%%%%%%%%%%%%%%%%%%%%%%%%%%%%%%%%%%%%%%%%%%%%%%
	
	\section{Spekulative physikalische Interpretation}
	
	Die bisherige Konstruktion ist rein mathematisch.
	
	Dennoch legen die numerischen Untersuchungen
	eine mögliche Interpretation nahe:
	
	\[
	E
	\leftrightarrow
	\text{Kern},
	\]
	
	\[
	B,C
	\leftrightarrow
	\text{Hülle},
	\]
	
	\[
	A
	\leftrightarrow
	\text{Bindung}.
	\]
	
	Ebenso erinnern
	
	\[
	1er,\;2er,\;3er,\;4er
	\]
	
	an eine hierarchische Teilchenklassifikation.
	
	Ob diese Analogien physikalische Bedeutung besitzen,
	bleibt offen.
	
	%%%%%%%%%%%%%%%%%%%%%%%%%%%%%%%%%%%%%%%%%%%%%%%%%%%%%%%%%%%%
	
	\section{Offene Probleme}
	
	Zu den wichtigsten offenen Fragen gehören:
	
	\begin{enumerate}
		\item Zusammenhang mit SU(2).
		\item Zusammenhang mit SU(3).
		\item Dynamische Interpretation der Eta-Phase.
		\item Zusammenhang mit der Zetafunktion.
		\item Beziehung zu Hurwitz-Gittern.
		\item Beziehung zur Leech-Struktur.
		\item Herleitung realer Teilchenmassen.
		\item Herleitung realer Ladungen.
	\end{enumerate}
	
	%%%%%%%%%%%%%%%%%%%%%%%%%%%%%%%%%%%%%%%%%%%%%%%%%%%%%%%%%%%%
	
	\section{Schlussbemerkung}
	
	Das Bamberger Modell III beschreibt natürliche Zahlen
	durch eine vierstufige Hierarchie
	
	\[
	V_4
	\rightarrow
	Q_8
	\rightarrow
	\eta_{24}
	\rightarrow
	A_5.
	\]
	
	Dabei liefert
	
	\[
	V_4
	\]
	
	die Restklasse,
	
	\[
	Q_8
	\]
	
	die Orientierung,
	
	\[
	\eta_{24}
	\]
	
	die Phase,
	
	und
	
	\[
	A_5
	\]
	
	die orbitale Geometrie.
	
	Ob diese Struktur lediglich eine mathematische Analogie
	oder Ausdruck eines tieferen Zusammenhangs ist,
	muss zukünftige Forschung zeigen.
	
	\section{Bamberger Modell III}
	
	\subsection{Leitidee}
	
	Das Bamberger Modell III beschreibt natürliche Zahlen als Träger
	einer mehrschichtigen geometrischen Struktur.
	
	Die Theorie basiert auf zwei Hierarchien.
	
	Die erste Hierarchie besteht aus den normierten Divisionsalgebren
	
	\[
	\mathbb R
	\longrightarrow
	\mathbb C
	\longrightarrow
	\mathbb H
	\longrightarrow
	\mathbb O.
	\]
	
	Die zweite Hierarchie besteht aus den sichtbaren endlichen
	Projektionsstrukturen
	
	\[
	V_4
	\longrightarrow
	Q_8
	\longrightarrow
	\eta_{24}
	\longrightarrow
	A_5.
	\]
	
	Dabei liefert die erste Hierarchie die algebraische Tiefenstruktur,
	während die zweite Hierarchie deren diskrete Manifestation beschreibt.
	
	\bigskip
	
	Die zentrale These lautet:
	
	\begin{quote}
	Natürliche Zahlen besitzen neben ihrer Primfaktorzerlegung
	eine Orientierungs-, Phasen- und Orbitstruktur.
	\end{quote}
	
	\bigskip
	
	%%%%%%%%%%%%%%%%%%%%%%%%%%%%%%%%%%%%%%%%%%%%%%%%%%%%%%%%%%%%
	
	\subsection{Die tetraedrische Spinstruktur}
	
	Die invertierbaren Restklassen modulo \(12\)
	
	\[
	(\mathbb Z/12\mathbb Z)^\times
	=
	\{1,5,7,11\}
	\]
	
	werden durch
	
	\[
	E,A,B,C
	\]
	
	bezeichnet.
	
	Die Zuordnung
	
	\[
	E=(1,1,1)
	\]
	
	\[
	A=(-1,-1,1)
	\]
	
	\[
	B=(1,-1,-1)
	\]
	
	\[
	C=(-1,1,-1)
	\]
	
	erzeugt ein reguläres Tetraeder.
	
	Es gilt
	
	\[
	E+A+B+C=0.
	\]
	
	Die vier Zustände bilden daher einen
	spinartigen Orientierungsraum.
	
	Der Spin ist im Modell eine lokale Eigenschaft.
	
	Wir schreiben
	
	\[
	S=s(H),
	\qquad
	H\in\{E,A,B,C\}.
	\]
	
	Der Spin benötigt keine Bewegung.
	
	Er existiert bereits auf einem einzelnen Zustand.
	
	%%%%%%%%%%%%%%%%%%%%%%%%%%%%%%%%%%%%%%%%%%%%%%%%%%%%%%%%%%%%
	
	\subsection{Quaternionische Chiralität}
	
	Die EABC-Struktur wird quaternionisch gehoben:
	
	\[
	E\mapsto1,
	\qquad
	A\mapsto i,
	\qquad
	B\mapsto j,
	\qquad
	C\mapsto k.
	\]
	
	Dann gilt
	
	\[
	BC=+A,
	\qquad
	CB=-A.
	\]
	
	Hier entsteht die erste Orientierung.
	
	Wir definieren
	
	\[
	A_R=BC,
	\qquad
	A_L=CB.
	\]
	
	Der Zustand \(A\) wird damit als
	zweikomponentiger Bindungsspinor interpretiert:
	
	\[
	\Psi_A=
	\begin{pmatrix}
	A_R\\
	A_L
	\end{pmatrix}.
	\]
	
	Die Chiralität des Modells sitzt somit
	nicht in den Zuständen \(B\) und \(C\),
	sondern im Bindungskanal \(A\).
	
	%%%%%%%%%%%%%%%%%%%%%%%%%%%%%%%%%%%%%%%%%%%%%%%%%%%%%%%%%%%%
	
	\subsection{Eta-24-Phasen}
	
	Die Dedekindsche Eta-Funktion
	
	\[
	\eta(\tau)
	=
	q^{1/24}
	\prod_{n\ge1}(1-q^n)
	\]
	
	führt auf eine diskrete Phasenstruktur
	
	\[
	\varepsilon
	\in
	\mathbb Z_{24}.
	\]
	
	Der Vorzeichenwechsel
	
	\[
	A_R
	\leftrightarrow
	A_L
	\]
	
	entspricht
	
	\[
	\varepsilon
	\rightarrow
	\varepsilon+12.
	\]
	
	Die Eta-Phase speichert damit die Orientierung
	des Spinors.
	
	%%%%%%%%%%%%%%%%%%%%%%%%%%%%%%%%%%%%%%%%%%%%%%%%%%%%%%%%%%%%
	
	\subsection{Bahndrehimpuls als zentrale Größe}
	
	Der wichtigste Schritt des Bamberger Modells III
	besteht in der strikten Trennung von
	
	\[
	S
	\]
	
	und
	
	\[
	L.
	\]
	
	Der Spin ist lokal.
	
	Der Bahndrehimpuls entsteht erst durch einen Umlauf.
	
	\bigskip
	
	Für einen geschlossenen Pfad
	
	\[
	P_1
	\rightarrow
	P_2
	\rightarrow
	\cdots
	\rightarrow
	P_m
	\rightarrow
	P_1
	\]
	
	definieren wir
	
	\[
	L(P)
	=
	\sum_{r=1}^{m}
	P_r
	\times
	P_{r+1}.
	\]
	
	Dies ist die diskrete Version der klassischen Beziehung
	
	\[
	\mathbf L
	=
	\mathbf r
	\times
	\mathbf p.
	\]
	
	\bigskip
	
	Die numerischen Untersuchungen zeigen:
	
	\begin{enumerate}
	\item Dreieckszyklen des Tetraeders erzeugen ein duales Tetraeder.
	\item Viererzyklen erzeugen die drei Raumachsen.
	\item Der Bahndrehimpuls ist daher keine lokale,
	sondern eine globale Größe.
	\end{enumerate}
	
	Damit gilt:
	
	\[
	\boxed{
	S=\text{Punkt}
	}
	\]
	
	\[
	\boxed{
	L=\text{geschlossener Umlauf}
	}
	\]
	
	Dies ist die zentrale strukturelle Trennung
	des Modells.
	
	%%%%%%%%%%%%%%%%%%%%%%%%%%%%%%%%%%%%%%%%%%%%%%%%%%%%%%%%%%%%
	
	\subsection{Ikosaedrische Orbitstruktur}
	
	Während der Spin auf dem Tetraeder lebt,
	entsteht der Bahndrehimpuls auf einer höheren Ebene.
	
	Diese wird durch die Ikosaedergruppe
	
	\[
	A_5
	\]
	
	beschrieben.
	
	Der Goldene Schnitt
	
	\[
	\varphi
	=
	\frac{1+\sqrt5}{2}
	\]
	
	erscheint hierbei nicht als numerologische Konstante,
	sondern als Eigenwertstruktur des Ikosaeders.
	
	Die Orbitgeometrie des Modells wird daher
	durch einen ikosaedrischen Bahnimpuls
	
	\[
	L_\varphi
	\]
	
	beschrieben.
	
	%%%%%%%%%%%%%%%%%%%%%%%%%%%%%%%%%%%%%%%%%%%%%%%%%%%%%%%%%%%%
	
	\subsection{Topologische Holonomie}
	
	Der Bahndrehimpuls allein genügt nicht.
	
	Zusätzlich besitzt jeder Umlauf
	eine topologische Phase.
	
	Für Übergangsoperatoren
	
	\[
	U_{ij}
	\]
	
	definieren wir die Holonomie
	
	\[
	\Gamma
	=
	\prod U_{ij}.
	\]
	
	Diese Größe ist die arithmetische Analogie zu
	
	\begin{itemize}
	\item Aharonov--Bohm-Phasen,
	\item Berry-Phasen,
	\item topologischen Phasen des Quanten-Hall-Effekts.
	\end{itemize}
	
	Die Eta-24-Struktur klassifiziert diese Holonomien.
	
	\[
	\Gamma
	\in
	\mathbb Z_{24}.
	\]
	
	Damit erhält die Theorie eine genuine
	topologische Komponente.
	
	%%%%%%%%%%%%%%%%%%%%%%%%%%%%%%%%%%%%%%%%%%%%%%%%%%%%%%%%%%%%
	
	\subsection{Morley-Kopplung}
	
	Morley erscheint im endgültigen Modell
	nicht als fundamentale Struktur.
	
	Morley wirkt vielmehr als Kopplungsoperator
	zwischen Spinraum und Orbitraum.
	
	Wir schreiben formal
	
	\[
	\mu_M
	=
	\langle
	\psi_{V_4},
	\psi_{A_5}
	\rangle.
	\]
	
	Die Morley-Kopplung misst somit den Überlapp
	zwischen tetraedrischen und ikosaedrischen Eigenräumen.
	
	%%%%%%%%%%%%%%%%%%%%%%%%%%%%%%%%%%%%%%%%%%%%%%%%%%%%%%%%%%%%
	
	\subsection{Gesamtdrehimpuls}
	
	Der Gesamtdrehimpuls des Modells lautet
	
	\[
	\boxed{
	J_{\rm BM}
	=
	S
	+
	\mu_M L_\varphi
	+
	\chi(\Gamma).
	}
	\]
	
	Dabei bezeichnet
	
	\[
	S
	\]
	
	den tetraedrischen Spin,
	
	\[
	L_\varphi
	\]
	
	den ikosaedrischen Bahndrehimpuls,
	
	\[
	\Gamma
	\]
	
	die topologische Holonomie,
	
	und
	
	\[
	\mu_M
	\]
	
	die Morley-Kopplung.
	
	\bigskip
	
	Die Dynamik des Modells entsteht damit primär
	nicht aus dem Spin,
	sondern aus der Wechselwirkung von
	
	\[
	L_\varphi
	\]
	
	und
	
	\[
	\Gamma.
	\]
	
	%%%%%%%%%%%%%%%%%%%%%%%%%%%%%%%%%%%%%%%%%%%%%%%%%%%%%%%%%%%%
	
	\subsection{Zusammenfassung}
	
	Die endgültige Architektur des Bamberger Modells III lautet
	
	\[
	\mathbb R
	\rightarrow
	\mathbb C
	\rightarrow
	\mathbb H
	\rightarrow
	\mathbb O
	\]
	
	\[
	\Downarrow
	\]
	
	\[
	V_4
	\rightarrow
	Q_8
	\rightarrow
	\eta_{24}
	\rightarrow
	A_5
	\]
	
	\[
	\Downarrow
	\]
	
	\[
	S
	\rightarrow
	L_\varphi
	\rightarrow
	\Gamma
	\]
	
	\[
	\Downarrow
	\]
	
	\[
	J_{\rm BM}.
	\]
	
	Der Spin beschreibt die lokale Orientierung.
	
	Der Bahndrehimpuls beschreibt die globale Umlaufstruktur.
	
	Die Holonomie beschreibt die topologische Phase.
	
	Der Gesamtdrehimpuls entsteht aus dem Zusammenspiel
	dieser drei Ebenen.
	
	\section{Kepler-, E8- und Leech-Füllungsschichten}
	
	Die bisherige Struktur des Bamberger Modells III beschreibt Spin,
	Bahndrehimpuls und Holonomie. Eine weitere notwendige Schicht ist die
	Füllungsdichte. Sie misst, ob eine geometrische Struktur lediglich
	orientiert ist oder ob sie eine maximal rigide Packung bildet.
	
	In drei Dimensionen liefert die Kepler-Packung die maximale Dichte
	kongruenter Kugeln:
	\[
	\delta_3=\frac{\pi}{\sqrt{18}}=\frac{\pi}{3\sqrt2}.
	\]
	Diese Konstante wird als dreidimensionale Füllungsgrenze des Modells
	interpretiert.
	
	In acht Dimensionen tritt das \(E_8\)-Gitter auf. Nach Viazovskas
	Lösung des Kugelpackungsproblems in Dimension \(8\) besitzt die optimale
	Packung die Dichte
	\[
	\delta_8=\frac{\pi^4}{384}.
	\]
	Damit erhält die oktonische Schicht des Modells eine konkrete
	Packungsgeometrie.
	
	In Dimension \(24\) erscheint das Leech-Gitter. Es steht in natürlicher
	Nähe zur Eta-\(24\)-Phase des Modells und ergänzt die
	\(24\)-periodische Spinorphase durch eine \(24\)-dimensionale
	Packungsrigidität.
	
	Wir erweitern daher die Massenformel zu
	\[
	m_{\rm BM}
	=
	\Lambda(\eta,\varphi)\,
	\rho_{\rm pack}
	\left(
	\Omega+\kappa\|S\|^2+\lambda\|L_\varphi\|^2
	+\theta\|\Gamma\|^2
	\right).
	\]
	Hierbei bezeichnet \(\rho_{\rm pack}\) eine geeignete
	Füllungsdichte, etwa \(\delta_3\), \(\delta_8\) oder eine
	Leech-artige \(24\)-dimensionale Dichte.
	
	\bibliographystyle{plain}
	
	\begin{thebibliography}{9}
		
		\bibitem{euler}
		Leonhard Euler.
		\newblock Introductio in Analysin Infinitorum.
		
		\bibitem{hamilton}
		William Rowan Hamilton.
		\newblock Lectures on Quaternions.
		
		\bibitem{hurwitz}
		Adolf Hurwitz.
		\newblock Vorlesungen über Zahlentheorie.
		
		\bibitem{dedekind}
		Richard Dedekind.
		\newblock Erläuterungen zu den modularen Funktionen.
		
		\bibitem{ramanujan}
		Srinivasa Ramanujan.
		\newblock Collected Papers.
		
		\bibitem{klein}
		Felix Klein.
		\newblock Vorlesungen über das Ikosaeder.
		
	\end{thebibliography}
	
\end{document}